\documentclass{article}
\usepackage[utf8]{inputenc}
\usepackage[T1]{fontenc}
\usepackage{geometry}
\geometry{
    a4paper,
    left=2.5cm,
    right=2.5cm,
    top=2.5cm,
    bottom=2.5cm
}
\usepackage{graphicx}
\usepackage{float}
\usepackage{hyperref}
\usepackage[english]{babel}
\usepackage{amsmath}
\usepackage{amssymb}
\usepackage{array}
\usepackage{booktabs}
\usepackage{caption}
\usepackage{subcaption}
\usepackage{enumitem}
\usepackage{fancyhdr}
\usepackage{titlesec}
\usepackage{url}
\usepackage{textcomp}
\usepackage{xcolor}
\usepackage{listings}
\usepackage{multirow}
\usepackage{longtable}
\usepackage{tabularx}
\usepackage{times}

% Set up page style
\pagestyle{fancy}
\fancyhf{}
\rfoot{\thepage}
\renewcommand{\headrulewidth}{0pt}

% Title formatting
\titleformat{\chapter}[display]{\normalfont\bfseries}{\chaptertitle\thechapter}{0pt}{\Huge}
\titlespacing*{\chapter}{0pt}{50pt}{40pt}

% Document info
\title{Template Document}
\author{Academic Support Agent}
\date{\today}

\begin{document}
\maketitle

%----------- Informations du rapport ---------

\logoentreprise{logos/logoGraphene.png}

\titre{Rapport du stage de fin d' tude} % Titre du fichier

\mention{Mention Sciences des donn es et de l'information} % Nom de la Mention
\trigrammemention{SDI} % Pour le bas de la page
\master{Master MSc IA} % Nom du master
\filiere{Fili re Management de projet et Transformation} % Nom de la fili re

\eleve{Pr nom Nom}

\dates{date d but - date fin}

% Informations tuteurs  coles
\tuteurecole{
    Mention : \textsc{Pr nom Nom} \\
    prenom.nom@centralesupelec.fr \\
    Fili re : \textsc{Pr nom Nom} \\ 
    pr nom.nom@gmail.com 
} 

\tuteurentreprise{
    \textsc{Pr nom Nom} \\
    pr nom.nom@entreprise.com 
}

%----------- Initialisation -------------------
        
\fairemarges %Afficher les marges
\fairepagedegarde %Cr er la page de garde

%----------- Abstract -------------------
\vspace*{\stretch{1}}
\begin{center}
	\begin{abstract}
        \lipsum[1-2]
    \end{abstract}
\end{center}
\vspace*{\stretch{1}}
\newpage

%------------ Table des mati res ----------------

\tabledematieres % Cr er la table de mati res

%------------ Corps du rapport ----------------


%------------ Introduction ----------------

\section{Introduction} 
\lipsum[3-4] % Effacer cette ligne et  crire le texte souhait 




\newpage
%------------ Related work ----------------

\section{Travaux connexes}
\lipsum[1-2] 





\newpage
%------------ Methodology ----------------

\section{M thodologie}
\lipsum[2-3]





\newpage
%------------ Application   l'analyse de texte ----------------

\section{Section 1}



%------------- Commandes utiles ----------------

\section{Quelques commandes}

Voici quelques commandes utiles :

%------ Pour ins rer et citer une image centralis e -----

\insererfigure{logos/logoCS.png}{3cm}{L gende de la figure}{Label de la figure}
% Le premier argument est le chemin pour la photo
% Le deuxi me est la hauteur de la photo
% Le troisi me la l gende
% Le quatri me le label

Ici, je cite l'image \ref{fig: Label de la figure}


%------- Pour ins rer et citer une  quation --------------

\begin{equation} \label{eq: exemple}
\rho + \Delta = 42
\end{equation}

L' quation \ref{eq: exemple} est cit  ici.
\end{document}